\section{Introduction}\label{sec:introduction}

Forecasting stock price movements has remained an important research
problem in financial analytics for many years. Among the various
approaches that have been proposed, Long Short-Term Memory (LSTM)
networks have become particularly popular because of their ability to
capture long-term patterns in sequential financial data. More recently,
Large Language Models (LLMs) have demonstrated promising capabilities in
analysing textual information obtained from financial news and market
reports.

Although previous studies have reported improvements through hybrid
approaches, the contribution of individual components remains unclear.
Improvements may originate from retrieval mechanisms, prompt design,
model architecture, or differences in the evaluation dataset.
Consequently, identifying the exact source of improvement becomes
difficult.

To address this limitation, the present study adopts a controlled
experimental framework in which four forecasting systems are evaluated
under identical conditions. The first model consists of a conventional
LSTM network that uses historical market data. The second model replaces
the LSTM with an LLM operating without external information. The third
model incorporates retrieval-based grounding, while the final model
combines the deterministic model and the grounded LLM through a
calibrated fusion strategy.

This approach makes it possible to evaluate the individual contributions
of language reasoning, retrieval, and model fusion while ensuring fair
comparison across all experiments.

\subsection{Research Question and Objectives}\label{sec:research-question}

This research aims to understand the individual contribution of
different components used in a hybrid stock market forecasting system.
The study investigates the following research question:

\textbf{Research Question:\\
}\emph{In five-day-ahead stock direction prediction, how much do
language reasoning, retrieval-based grounding, and calibrated fusion
with a deterministic sequence model individually contribute when all
models are evaluated on the same dataset, compared against a market
regime-aware baseline, and assessed using both prediction accuracy and
response quality?}

To answer this question, the study defines the following research
objectives:

\begin{itemize}
\item
  \textbf{O1.} Develop a deterministic machine learning model by
  comparing five candidate algorithms using the same walk-forward
  validation strategy and selecting the best-performing architecture
  based on experimental results rather than assumptions.
\item
  \textbf{O2.} Design two Large Language Model (LLM) systems that use
  the same model, prompt, temperature, and output format, with the only
  difference being the inclusion of retrieved external evidence.
\item
  \textbf{O3.} Create a hybrid forecasting model by combining the
  deterministic model and the grounded LLM through a calibrated fusion
  method, ensuring that the hybrid model provides meaningful
  contributions beyond simply favouring the stronger individual model.
\item
  \textbf{O4.} Perform a fully reproducible hallucination analysis using
  automated evaluation methods without relying on human reviewers or
  language-model-based judges.
\item
  \textbf{O5.} Evaluate the consistency of each forecasting approach by
  repeating identical inputs multiple times and measuring prediction
  stability instead of assuming deterministic behaviour from a
  zero-temperature setting.
\item
  \textbf{O6.} Assess the practical usefulness of each system by
  evaluating not only prediction accuracy but also the quality of the
  generated response, including its explanation, supporting evidence,
  reasoning, and investment guidance.
\item
  \textbf{O7.} Determine the maximum achievable prediction accuracy for
  the proposed forecasting task, allowing the reported results to be
  interpreted against a realistic upper performance bound rather than a
  theoretical baseline.
\end{itemize}

\subsection{Contribution and Structure}\label{sec:contribution}

The primary contribution of this research is not the development of a
stock market prediction model with the highest accuracy. Instead, it
introduces a systematic and reproducible evaluation framework that
measures the individual contribution of each component within a hybrid
forecasting system. By evaluating language reasoning, retrieval-based
grounding, and model fusion independently under identical conditions,
the proposed framework provides a clear understanding of the benefits
and limitations of each engineering layer.

An important aspect of this study is that both positive and negative
findings are reported with equal importance. The experimental results
show that retrieval improves the quality of generated responses but also
increases numerical unit errors. Similarly, the hybrid model reduces the
variation in prediction confidence compared with the standalone language
model, although it does not completely eliminate inconsistent
predictions. In addition, none of the observed improvements in
forecasting accuracy achieved statistical significance for the evaluated
dataset. Reporting these findings provides a balanced assessment of the
proposed approach and highlights both its strengths and limitations.

The remainder of this paper is organized as follows. \textbf{Section 2}
reviews the existing literature and identifies the research gap
addressed by this study. \textbf{Section 3} explains the research
methodology, including the dataset, evaluation protocol, and
experimental design. \textbf{Section 4} presents the system architecture
and describes the proposed fusion algorithm. \textbf{Section 5}
discusses the implementation details and the developed software
components. \textbf{Section 6} reports the experimental results,
evaluates the proposed models through five experiments, and discusses
the key findings. Finally, \textbf{Section 7} summarizes the conclusions
of the study and outlines possible directions for future research.
